\section{第11章：光学（11.5练习题）}
\sourceproblem{选择1}{199}{205}
\problemtext 两个独立白光源的光束重叠，判断没有稳定干涉图样的主要原因。
\workings 时间平均强度为 $I_1+I_2+2\operatorname{Re}\langle E_1 E_2^*\rangle$（统一归一化）。独立普通光源的互相关平均为零，故无稳定干涉条纹，选D。
\crosscheck 同一白光源分束在足够小光程差处仍可干涉，白光多波长不是充分原因。
\sourceproblem{选择2}{199}{205}
\problemtext 在棱镜色散、彩虹、水珠透亮和油膜彩色花纹中识别薄膜干涉。
\workings 油膜上下表面的反射光相干叠加，不同波长对应不同增强条件，产生彩色花纹，选D。
\crosscheck 彩虹的主要几何光学机制是折射、色散及内部反射；不把一切彩色现象归为干涉。
\sourceproblem{选择3}{199}{205}
\problemtext 白光双缝的两缝分别加理想红、蓝滤片，判断干涉图样。
\workings 两透射频谱不重叠，探测时间远长于差频周期，交叉项平均为零，故无双缝稳定干涉，选D。
\crosscheck 若实际滤片存在共同透射频段，可能有微弱残余干涉；每个单缝自身仍可衍射。
\sourceproblem{选择4}{199}{205}
\problemtext 双缝一缝略变窄，比较条纹间距及极小值。
\workings 缝中心距d不变、近轴模型中间距 $\Delta x=\lambda D/d$ 不变；到达强度不等时 $I_{\min}=(\sqrt{I_1}-\sqrt{I_2})^2>0$，选C。
\crosscheck 假设缝中心不移动且局部单缝包络变化可忽略；有限缝宽还影响衍射包络。
\sourceproblem{选择5}{199}{205}
\problemtext 原双缝P点为亮纹，遮挡一缝并在两缝中垂面设置高折射率介质反射镜，判断P点。
\workings 理想洛埃镜模型中虚源与被遮缝对称位置重合；原几何光程差对应整数波长，反射相位额外增加π，故P为干涉极小，选B。
\crosscheck 这里暗纹指极小，不必为零；实际镜的偏振、入射角及复反射系数会影响反射相位，教材采用理想半波损失模型。
\sourceproblem{选择6}{199}{205}
\problemtext 波长600 nm、缝中心距2 mm、屏距3 m，求相邻明纹间距（题尾续PDF206）。
\workings 近轴杨氏双缝 $\Delta x=\lambda D/d=0.90$ mm，选B。
\crosscheck 小角度条件下明纹近似等间隔；远离中心需使用精确光程差。
\sourceproblem{选择7}{200}{206}
\problemtext 真空波长 $\lambda$ 的光在折射率n介质内沿AB传播，未折返传播相差 $3\pi$，求光程。
\workings $\Delta\phi=(2\pi/\lambda)\mathcal L$，故 $\mathcal L=1.5\lambda$，选A。
\crosscheck 几何长度为 $1.5\lambda/n$；若只给模 $2\pi$ 的相差，则光程还存在整数波长的不确定性。
\sourceproblem{选择8}{200}{206}
\problemtext 空气、折射率1.3油膜和1.50玻璃构成薄膜，正入射500 nm光，求透射增强的最小正厚度。
\workings 两次反射均由低折射率到高折射率，反射相对附加相位为零。反射极小满足 $2nd=\lambda/2$（最低阶），无吸收时透射极大，故 $d=\lambda/(4n)\approx96.15$ nm，选D。
\crosscheck 反射极小不一定零反射，取决于两界面反射幅度；此处正厚度最小。
\sourceproblem{选择9}{200}{206}
\problemtext 空气中1.60玻璃表面镀1.38的MgF2膜，500 nm正入射，求最小增透膜厚度。
\workings 两界面反射相位均反转；令往返传播相位为π，$d=500/(4\cdot1.38)\approx90.58$ nm，选B。
\crosscheck 完全消反射还需 $n_2^2=n_1n_3$；给定膜折射率不满足此条件，故仅使反射尽可能小。
\sourceproblem{选择10}{200}{206}
\problemtext 牛顿环夹层空气抽成真空，比较环半径。
\workings 对接触装置反射暗环 $r_m^2\simeq m\lambda R/n$，固定级次和真空波长时n减小使半径增大，选A。
\crosscheck 明环也具有同样的 $n^{-1/2}$ 缩放；空气折射率变化很小，实际半径变化亦小。
\sourceproblem{选择11}{200}{206}
\problemtext 图示上透镜1.52、夹层1.62，底板左半1.75、右半1.52，判断接触处反射条纹明暗。
\workings 第一反射低到高相位反转；左半下界面亦低到高，两反射同附加相位，零厚度为极大。右半下界面高到低，无反转，相差π，零厚度为极小，选D。
\crosscheck 明暗为相对极值；不同反射幅值时暗侧不必为零。已按原图核对三个区域折射率。
\sourceproblem{选择12}{200}{206}
\problemtext 牛顿环平凸透镜缓慢上移，判断环运动和中心亮暗。
\workings 空气层 $d(r)\simeq h+r^2/(2R)$，固定暗环级次满足 $r_m^2=m\lambda R-2Rh$，故h增大使环向中心收缩；中心随 $2h/\lambda$ 变化明暗交替，选B。
\crosscheck 固定相邻级次平方半径差不变；中心环消失后应重新标记可见环，不能笼统说间隔变小。
\sourceproblem{选择13}{200}{206}
\problemtext 1.35油滴在1.60玻璃上形成中央厚1微米的薄膜，600 nm正入射，判断反射边缘和中心。
\workings 上下反射均低到高，附加相位相消。边缘 $d=0$ 为极大；中心 $2nd/\lambda=2(1.35)(1000)/600=4.5$，为极小，选A。
\crosscheck 半整数往返光程对应反射极小；原题膜形只决定条纹位置，不改变局部厚度条件。
\sourceproblem{选择14}{201}{207}
\problemtext 单缝衍射图中P为中央亮纹一侧第一暗纹，求图示BC的光程差。
\workings 图中C为从A向B出射方向作垂线所得投影，$BC=b\sin\theta$。第一暗纹满足 $b\sin\theta=\lambda$，故BC为 $\lambda$，选B。
\crosscheck 半波带宽对应 $\lambda/2$，第一暗纹由两个半波带相消；BC为全缝端点差。
\sourceproblem{选择15}{201}{207}
\problemtext 在几何影子、树荫亮斑、金属丝阴影中央亮线和透镜光斑中识别衍射现象。
\workings 金属丝两侧绕射光在几何阴影内叠加可形成中央亮线，选C。
\crosscheck 其余现象主要可由几何光学解释，但真实光学系统也有衍射修正。
\sourceproblem{选择16}{201}{207}
\problemtext 波长500 nm、缝宽0.25 mm，焦平面两侧第三暗纹相距12 mm，求焦距。
\workings $b\sin\theta_3=3\lambda$，$\sin\theta_3=0.006$，$2f\tan\theta_3=0.012$ m，故 $f\simeq1.0$ m，选B。
\crosscheck 用精确三角关系得约0.999982 m；小角近似与题目精度一致。
\sourceproblem{选择17}{201}{207}
\problemtext 夫琅禾费单缝装置中，平面波照明，单缝沿光轴向透镜移动，观察焦面条纹。
\workings 理想透镜将同一方向的平面波聚焦到相同焦面位置；轴向平移给角谱增加相位因子而不改变其模，故图样不变，选C。
\crosscheck 假设透镜口径足够、屏位于焦面且无截光；有限口径或非平面照明可能改变强度分布。
\sourceproblem{选择18}{201}{207}
\problemtext 单缝夫琅禾费衍射中减小缝宽，比较中央明纹宽度。
\workings 第一暗纹 $\sin\theta_1=\lambda/b$；近轴焦面中央宽度 $2f\lambda/b$，b减小使宽度增大，选B。
\crosscheck 公式须有第一暗纹（$b\ge\lambda$）；当缝极窄时不再有有限角度的第一暗纹。
\sourceproblem{选择19}{201}{207}
\problemtext 第一级侧明纹的衍射角0.001 rad，估计缝宽。
\workings 教材半波带近似 $b\sin\theta\simeq3\lambda/2$，给 $b\simeq1500\lambda$，选D。严格均匀单缝强度 $I/I(0)=(\sin\beta/\beta)^2$，$\beta=\pi b\sin\theta/\lambda$；首侧极大满足 $\tan\beta=\beta$，根 $\beta\simeq4.493409$，给 $b\simeq1430.3\lambda$。
\crosscheck 1500是近似而非精确极大位置；原选项中D符合教材近似。
\sourceproblem{选择20}{201}{207}
\problemtext 单缝在30度方向恰分为3个半波带，求缝宽。
\workings $b\sin30^\circ=3\lambda/2$，故 $b=3\lambda$，选D。
\crosscheck 半波带数由全缝光程差除以 $\lambda/2$ 得到。
\sourceproblem{选择21}{201}{207}
\problemtext 波长550 nm、圆孔直径127 cm望远镜，按瑞利判据求最小分辨角。
\workings $\theta_{\min}\simeq1.22\lambda/D=1.22(550\times10^{-9})/1.27\simeq5.28\times10^{-7}$ rad，选D。
\crosscheck D是孔径直径，非半径；忽略像差、大气扰动等。
\sourceproblem{选择22}{201}{207}
\problemtext 白光正入射光栅，同一级可见光谱中哪种颜色偏离中心最远（选项续PDF208）。
\workings $d\sin\theta=m\lambda$，固定非零级次时角度绝对值随波长增大，红光最远，选D。
\crosscheck 仅比较实际存在的同级谱线；若 $|m|\lambda>d$ 则该谱线不存在。
\sourceproblem{选择23}{202}{208}
\problemtext 光栅常数 $9\times10^{-6}$ m，第三级主极大 $\sin\theta=0.2$，求波长。
\workings $\lambda=d\sin\theta/3=600$ nm，选C。
\crosscheck 此为主极大位置条件；是否缺级还取决于透光缝宽。
\sourceproblem{选择24}{202}{208}
\problemtext 光栅周期 $d=2b$，b为透光缝宽，判断缺级。
\workings 主极大 $d\sin\theta=m\lambda$，单缝零点 $b\sin\theta=p\lambda$（非零整数p），联立 $m=2p$，故非零偶数级缺失，选C。
\crosscheck 零级不缺失；只讨论满足传播角条件的级次。
\sourceproblem{选择25}{202}{208}
\problemtext 旋转理想线偏振片时透射光强变化但无全暗，判断入射偏振状态。
\workings 教材若仅考虑自然光、线偏振光及二者非相干混合，则变化而不消光说明有线偏振成分及非零背景，选B。但完全椭圆偏振光也可产生 $I(\theta)=a^2\cos^2\theta+b^2\sin^2\theta$（两正交分量相差π/2，$a,b$非零且不等）。
\crosscheck 一般偏振态中仅凭旋转线分析器不能排除完全椭圆偏振；若不限定教材三类模型，则选D更严谨。
\sourceproblem{选择26}{202}{208}
\problemtext 两偏振片轴夹角从60度变为45度，入射强度不变，求透射强度比。
\workings 若第一片后强度固定，马吕斯定律给 $I_{45}/I_{60}=\cos^245^\circ/\cos^260^\circ=2$，选A。
\crosscheck 入射为自然光时任意转动第一片仍给固定首片透射强度；若入射为偏振光且第一片方向也变动，仅给入射总强度不够确定该比。
\sourceproblem{选择27}{202}{208}
\problemtext 自然光以60度入射界面，反射光完全线偏振，判断折射光及折射角。
\workings 对普通透明各向同性非磁介质，布儒斯特角满足 $i_B+r=90^\circ$，故折射角30度。透射光通常仍有两个正交偏振分量，为部分偏振光，选D。
\crosscheck 由 $n_2/n_1=\tan60^\circ$ 确定相对折射率而非要求入射介质为真空。
\sourceproblem{选择28}{202}{208}
\problemtext 自然光从空气以布儒斯特角入射玻璃，判断反射光偏振方向。
\workings 此角下p偏振的菲涅耳反射系数为零，剩余s偏振，电场振动垂直入射面，完全线偏振，选C。
\crosscheck 对普通透明各向同性非磁介质成立；透射光一般并非完全偏振。
\sourceproblem{选择29}{202}{208}
\problemtext 某光滑介质表面在58度入射时反射光为线偏振，求折射率。
\workings 若从空气入射且入射光含两偏振分量、完全偏振由布儒斯特条件产生，$n/n_{\rm air}=\tan58^\circ$，取空气折射率1，选C。
\crosscheck 原题未明确入射介质及入射偏振态，严格只能在上述条件下推出相对折射率；已线偏振的入射光不能据此判定布儒斯特角。
\sourceproblem{选择30}{202}{208}
\problemtext 自然光以起偏角入射平板玻璃，选择入射角和折射角关系。
\workings 布儒斯特条件 $i_B+r=90^\circ$，选C；此时反射光线与折射光线互相垂直。
\crosscheck 该关系指第一界面折射角；平行平板最终出射光与原入射光平行，不应混淆。

\sourceproblem{填空1}{202}{208}
\problemtext 等光程照明双缝，空气中P为第三级明纹，浸入液体后P为第四级明纹（题尾与图在PDF209），求原光程差及液体折射率。
\workings 取图中P靠近S1，定义 $\Delta=r_2-r_1>0$，空气近似真空时 $\Delta=3\lambda$。浸液后光程差 $n\Delta=4\lambda$，故 $n=4/3$。
\crosscheck 假定几何不变且光源频率不变；若计空气折射率，所得为液体相对空气折射率。
\sourceproblem{填空2}{203}{209}
\problemtext 波长562.5 nm、屏距1.2 m、明纹间距1.5 mm，求双缝间距。
\workings 近轴杨氏公式 $d=\lambda D/\Delta x=4.50\times10^{-4}$ m，即0.450 mm。
\crosscheck $\Delta x/D=0.00125$，小角近似合理。
\sourceproblem{填空3}{203}{209}
\problemtext 双缝S1覆盖厚e、折射率n的云母片，求中央明纹移动方向及原中心光程差。
\workings 空气近似真空，S1路新增光程 $(n-1)e$。原中心两几何路程相等，故S1路减S2路的光程差为 $(n-1)e$。零级需S1几何路程相对缩短，故向S1侧移动；近轴位移大小 $D(n-1)e/d$。
\crosscheck 假定光近垂直穿膜且 $n>1$；若反取差的定义，符号反号。
\sourceproblem{填空4}{203}{209}
\problemtext 厚e、折射率 $n_2$ 的平行薄膜夹在 $n_1,n_3$ 两介质之间，正入射，比较两种折射率排序的反射相消条件。
\workings 真空波长为 $\lambda$。若 $n_1<n_2>n_3$，两反射有一次相对π相变，极小条件为 $2n_2e=k\lambda$；若 $n_1<n_2<n_3$，两次均反相，极小条件为 $2n_2e=(k-1/2)\lambda$（$k=1,2,\ldots$）。
\crosscheck 原题称反射消失，但严格消光还须反射振幅匹配。对无吸收单层膜，前一种整数半波厚度需 $n_1=n_3$，后一种奇数四分之一波厚度需 $n_2^2=n_1n_3$；否则只是反射极小。
\sourceproblem{填空5}{203}{209}
\problemtext 波长600 nm的反射牛顿环，第四个非中央暗环的空气膜厚度。
\workings 接触式装置一处相对半波损失，暗环 $2e=k\lambda$，中央为 $k=0$。第四环 $e=4\lambda/2=1200$ nm。
\crosscheck 空气折射率近似1；不把中央暗斑算作第一暗环。
\sourceproblem{填空6}{203}{209}
\problemtext 空气中劈形透明膜，劈角 $10^{-4}$ rad，波长700 nm，明纹间距0.25 cm，求折射率。
\workings 相邻同类条纹厚差 $\lambda/(2n)$，小角下 $\Delta e=b\tan\theta\simeq b\theta$，所以 $n=\lambda/(2b\tan\theta)\simeq1.40$。
\crosscheck 0.25 cm等于0.0025 m；角度必须用弧度。
\sourceproblem{填空7}{203}{209}
\problemtext 空气劈尖角度减小，判断条纹疏密变化。
\workings 间距 $b\simeq\lambda/(2\theta)$ 增大，单位长度条纹数 $1/b$ 减小，条纹变疏。
\crosscheck 原书使用疏密一词，明确区分间距增大与密度减小。
\sourceproblem{填空8}{203}{209}
\problemtext 波长600 nm牛顿环第二与第五明环的膜厚差。
\workings 相邻明环厚度差为 $\lambda/2$，相隔3级，故 $\Delta e=3\lambda/2=900$ nm。
\crosscheck 取差的绝对值；与明环编号起点无关。
\sourceproblem{填空9}{203}{209}
\problemtext 接触式牛顿环用600和500 nm光，求两种光第五个反射明环的膜厚差。
\workings 从中心向外第m个明环 $e_m=(2m-1)\lambda/4$，故第五明环厚差 $9(600-500)/4=225$ nm。
\crosscheck 中央为暗斑，第一明环厚度 $\lambda/4$。
\sourceproblem{填空10}{203}{209}
\problemtext 单缝第三暗纹的半波带数，缝宽减半后原第三暗纹位置属于何种条纹。
\workings 原有 $b\sin\theta=3\lambda$，分成6个半波带。缝宽减半后 $\beta=\pi b'\sin\theta/\lambda=3\pi/2$，强度非零，为教材半波带近似的第一侧明纹。
\crosscheck 严格极大满足 $\tan\beta=\beta$，第一侧峰在4.4934而非 $3\pi/2$；原位置并非精确峰，归一化强度为 $4/(9\pi^2)$，是亮区而非暗纹。
\sourceproblem{填空11}{203}{209}
\problemtext 钠黄光589 nm的中央明纹宽4.0 mm，求442 nm光在同装置中的中央明纹宽。
\workings 小角下 $w=2f\lambda/b$，故 $w'=4.0(442/589)\simeq3.00$ mm。
\crosscheck 固定焦距、缝宽，波长较短时中央纹较窄。
\sourceproblem{填空12}{203}{209}
\problemtext 同一波长的单缝衍射，缝越窄中央明纹如何变化（题尾PDF210）。
\workings 第一暗纹 $\sin\theta_1=\lambda/b$，b减小使中央明纹变宽。
\crosscheck 若 $b<\lambda$ 则不存在有限角度第一暗纹，常规宽度定义不再适用。
\sourceproblem{填空13}{204}{210}
\problemtext 30度衍射方向对应4个半波带，求缝宽。
\workings $b\sin30^\circ=4\lambda/2=2\lambda$，所以 $b=4\lambda$。
\crosscheck 四个半波带成对相消，对应第二暗纹。
\sourceproblem{填空14}{204}{210}
\problemtext 单缝中央明纹旁第二暗纹对应半波带数。
\workings $b\sin\theta=2\lambda$，每半波带光程差 $\lambda/2$，故数目4。
\crosscheck 一般第m暗纹对应2m个半波带。
\sourceproblem{填空15}{204}{210}
\problemtext 瞳孔直径5 mm、波长400 nm，距地面点光源200 km，按衍射极限求最小分辨间距。
\workings 圆孔瑞利角 $\theta=1.22\lambda/D=9.76\times10^{-5}$ rad。两源横向距离 $s\simeq L\theta=19.52$ m。
\crosscheck 仅计理想圆孔衍射；实际人眼还受像差和视网膜采样限制，距离取垂直视线的分量。
\sourceproblem{填空16}{204}{210}
\problemtext 每毫米800线的光栅，一级谱线角30度，求波长。
\workings 周期 $d=10^{-3}/800=1.25\times10^{-6}$ m，$\lambda=d\sin30^\circ=625$ nm。
\crosscheck 刻线密度是周期的倒数。
\sourceproblem{填空17}{204}{210}
\problemtext 光栅透光缝宽等于不透光间隔，求可能出现的谱级。
\workings $d=2b$，光栅主极大与单缝零点重合时 $m=2p$（非零p），故零级及 $\pm1,\pm3,\pm5,\ldots$ 奇数级可能出现。
\crosscheck 还必须满足 $|m|\lambda\le d$；零级不属于缺级。
\sourceproblem{填空18}{204}{210}
\problemtext 波长500 nm，光栅常数 $10^{-4}$ cm，求一级主极大角。
\workings $d=10^{-6}$ m，$\sin\theta=\lambda/d=1/2$，故 $\theta=30^\circ$（两侧为 $\pm30^\circ$）。
\crosscheck cm换算为m时指数再减2。
\sourceproblem{填空19}{204}{210}
\problemtext 固定波长，提高光栅单位长度缝数，判断光栅常数和相邻明纹间距。
\workings $d=1/\eta$ 随刻线密度 $\eta$ 增大而减小，近轴主极大角间隔 $\Delta\theta\simeq\lambda/d$ 增大，焦面间距亦增大。
\crosscheck 较大角度时谱线并非等间距；固定可存在的同一对级次比较，且不讨论缺级。
\sourceproblem{填空20}{204}{210}
\problemtext 光栅周期、缝宽和波长固定，增加有效照明缝数N，判断同级谱线宽度。
\workings 主极大两侧最近零点的近似角半宽为 $\lambda/(Nd\cos\theta_m)$，故N增大使谱线变窄。
\crosscheck 主极大位置不变，峰值强度在各缝照明固定时按 $N^2$ 增长；若该级缺级则无此峰。
\sourceproblem{填空21}{204}{210}
\problemtext 线偏振光经最少块理想偏振片使方向转90度，求片数及最大透射比例。
\workings 单片正交时透射为零；两片可实现非零透射。中间轴与初始方向夹角 $\alpha$，末片正交，则 $I/I_0=\cos^2\alpha\sin^2\alpha=\sin^22\alpha/4$，在45度时最大为1/4。
\crosscheck 最少两片；若允许更多片，分步转动可提高透射率，故1/4是两片条件下最大值。
\sourceproblem{填空22}{204}{210}
\problemtext 自然光依次穿两理想偏振片，轴夹45度，末端光强I，求第二片入射光强。
\workings 马吕斯定律 $I=I_1\cos^245^\circ=I_1/2$，故 $I_1=2I$。
\crosscheck 原自然光强为4I，题问的不是这一数值。
\sourceproblem{填空23}{204}{210}
\problemtext 自然光强 $I_0$ 和线偏振光强I非相干混合，旋转分析器后最大光强为最小的5倍，求比例。
\workings $I_{\max}=I_0/2+I$，$I_{\min}=I_0/2$。联立得 $I=2I_0$，所以 $I_0:I=1:2$。
\crosscheck 自然光经任何理想线偏振片均透射一半；非相干分量强度直接相加。
\sourceproblem{填空24}{204}{210}
\problemtext 自然光强 $I_0$ 依次过三理想偏振片，相邻轴夹45度，求末端强度。
\workings $I=(I_0/2)\cos^245^\circ\cos^245^\circ=I_0/8$。
\crosscheck 第一片因自然光减半，后两片各按马吕斯定律减半。
\sourceproblem{填空25}{204}{210}
\problemtext 正交的P1、P3间放匀角速旋转的P2，初始P2与P1平行，自然入射强度 $I_0$，求P1、P3输出。
\workings $I_1=I_0/2$，P2轴角 $\omega t$，$I_3=(I_0/2)\cos^2\omega t\sin^2\omega t=(I_0/8)\sin^2(2\omega t)$。
\crosscheck t=0输出零，P2在45度时最大 $I_0/8$，强度周期为 $\pi/(2|\omega|)$。
\sourceproblem{填空26}{204}{210}
\problemtext 两理想偏振片透振方向夹30度，自然入射强度 $I_0$，求透射强度。
\workings $I=(I_0/2)\cos^230^\circ=3I_0/8$。
\crosscheck 不可漏掉起偏器的1/2因子。
\sourceproblem{填空27}{204}{210}
\problemtext 自然光以60度射向玻璃，反射光完全线偏振，求折射角及玻璃折射率。
\workings 在空气入射、普通透明非磁介质模型下，$r=90^\circ-60^\circ=30^\circ$，$n/n_{\rm air}=\tan60^\circ=\sqrt3$，取空气近似1则玻璃折射率 $\sqrt3$。
\crosscheck 由斯涅尔定律 $\sin60^\circ/\sin30^\circ=\sqrt3$ 复核。
\sourceproblem{填空28}{204}{210}
\problemtext 折射率1.65的火石玻璃浸在1.33的水中，求水侧布儒斯特角。
\workings $i_B=\arctan(n_{\rm glass}/n_{\rm water})=\arctan(1.65/1.33)\simeq51.13^\circ$。
\crosscheck 此处相对折射率不是1.65，须计水的折射率。
\sourceproblem{填空29}{204}{210}
\problemtext 某介质到空气的全反射临界角45度，求从空气射向此介质的布儒斯特角。
\workings $\sin45^\circ=n_{\rm air}/n$，故 $n/n_{\rm air}=\sqrt2$，$i_B=\arctan\sqrt2\simeq54.74^\circ$。
\crosscheck 临界角属于高折射率向低折射率入射；布儒斯特角按相反传播方向计算。
\sourceproblem{填空30}{205}{211}
\problemtext 自然光以布儒斯特角入射透明介质分界面，求反射光偏振类型及电场振动方向。
\workings 对普通各向同性透明非磁介质，p分量反射系数为零，反射光为完全线偏振光，电场垂直入射面。
\crosscheck 若两介质折射率完全相等，则无界面反射光，谈不上其偏振方向。
\sourceproblem{计算1}{205}{211}
\problemtext 双缝距0.3 mm、波长600 nm、屏距0.5 m，求第二及第三暗纹的距离。
\workings 近轴相邻暗纹间距 $\Delta x=\lambda D/d=1.0$ mm，第二、第三暗纹在同侧时相距1.0 mm。若在异侧，距中心分别 $1.5\Delta x,2.5\Delta x$，总距4.0 mm。
\crosscheck 原题未明确两纹所在侧，通常按同侧理解；两种几何约定均列出。
\sourceproblem{计算2}{205}{211}
\problemtext 波长550 nm、双缝距 $2\times10^{-4}$ m、屏距2 m，求两侧第10级明纹间距；一缝盖厚6.6微米、折射率1.58玻璃后，求零级移至原明纹的级次。
\workings 近轴 $\Delta x=\lambda D/d=5.5$ mm，两侧第10级相距 $20\Delta x=110$ mm。覆盖新增光程 $(n-1)e=3.828\times10^{-6}$ m，位移为 $3.828/0.550=6.96$ 个条纹间距，约原第7级明纹处，朝覆盖缝。
\crosscheck 原题小数给出6.96而非精确7级；不强行写成整数。近轴位移约38.28 mm。
\sourceproblem{计算3}{205}{211}
\problemtext S1、S2分别覆盖折射率1.4和1.7的等厚玻璃，零级移至原第五级明纹，波长480 nm，求膜厚及原中心级次。
\workings 设厚度d，新增S2路相对光程为 $(1.7-1.4)d=0.3d=5\lambda$，故 $d=8.0\times10^{-6}$ m。零级向S2侧移；原中心光程差 $5\lambda$，变为第五级明纹（有向级次取决于定义）。
\crosscheck 两片都替换等厚空气，空气项相消；必须用折射率差而非单片折射率。
\sourceproblem{计算4}{205}{211}
\problemtext 双缝距0.3 mm、屏距1.20 m，两侧第五暗纹间距22.78 mm，求波长。
\workings 第五暗纹距中心为 $4.5\lambda D/d$，两侧间距 $9\lambda D/d$，所以
\[
\lambda=\frac{0.02278(0.0003)}{9(1.20)}
\simeq6.328\times10^{-7}\ \mathrm m=632.8\ \mathrm{nm}.
\]
\crosscheck 高清原页数字为22.78 mm，非22.87 mm；近轴半角约0.00949 rad，近似合理。
